% Part: lambda-calculus
% Chapter: church-rosser
% Section: parallel-beta-reduction

\documentclass[../../../include/open-logic-section]{subfiles}

\begin{document}

\olfileid{lam}{cr}{pb}

\olsection{সমান্তরাল $\beta$-হ্রাস}

আমরা \emph{সমান্তরাল $\beta$-হ্রাস}-এর ধারণা প্রবর্তন করব এবং প্রমাণ
করব যে এটি চার্চ--রসার ধর্ম পূরণ করে।

\begin{defn}[সমান্তরাল $\beta$-হ্রাস, $\bredpar$] \ollabel{defn:bredpar}
  পদগুলির সমান্তরাল হ্রাস ($\bredpar$) আরোহমূলকভাবে নিচের মতো সংজ্ঞায়িত:
  \begin{enumerate}
    \item \ollabel{defn:bredpar1} $x \bredpar x$।
    \item \ollabel{defn:bredpar2} $N \bredpar N'$ হলে $\lambd[x][N] \bredpar
      \lambd[x][N']$।
% BN-SRC-293 TARGET-CORRECTION-BEGIN
\par\smallskip\noindent\textnormal{\small\textbf{উৎস-সংশোধন (BN-SRC-293):} হিমায়িত দ্বিতীয় বিধির পূর্বধারণা $N\xrightarrow{\beta}N'$ সাধারণ $\beta$-হ্রাস ব্যবহার করে, যদিও সংজ্ঞাটি সমান্তরাল $\beta$-হ্রাসের আরোহমূলক বিধি দিচ্ছে এবং পরের দুই প্রমাণ একই বিধিকে $N\bredpar N'$ পূর্বধারণাসহ ব্যবহার করে। বাংলা বিধিতে তাই সংজ্ঞায়িত সম্বন্ধ~$\bredpar$ পুনরুদ্ধার করা হয়েছে।} % BN-SRC-293 TARGET-CORRECTION-END
    \item \ollabel{defn:bredpar3} $P \bredpar P'$ এবং $Q \bredpar Q'$ হলে $PQ \bredpar
      P'Q'$।
    \item \ollabel{defn:bredpar4} $N \bredpar N'$ এবং $Q \bredpar Q'$ হলে
      $(\lambd[x][N])Q \bredpar \Subst{N'}{Q'}{x}$।
  \end{enumerate}
\end{defn}

সমান্তরাল $\beta$-হ্রাসে কোনো পদের যেকোনো সংখ্যক রিডেক্সকে এক ধাপে
হ্রাস করা যায়। এটি $\beta$-হ্রাস থেকে এই অর্থে আলাদা যে, শুধু মূল
পদে থাকা রিডেক্সগুলিই সংকুচিত করা যায়; সমান্তরাল $\beta$-হ্রাস থেকে
নতুন যে রিডেক্স তৈরি হয়, তাকে একই ধাপে সংকুচিত করা যায় না। যেমন,
$(\lambd[f][fx])(\lambd[y][y])$ পদটিকে সমান্তরাল $\beta$-হ্রাস করে
শুধু পদটি নিজে অথবা $(\lambd[y][y])x$ পাওয়া যায়, আর এক ধাপেই~$x$
পাওয়া যায় না। অথচ সাধারণ $\beta$-হ্রাসে পদটি~$x$-এ হ্রাস পায়;
কারণ এই রিডেক্সটি সমান্তরাল $\beta$-হ্রাসের প্রথম ধাপের পরেই তৈরি হয়।
তবে সমান্তরাল $\beta$-হ্রাসের দ্বিতীয় ধাপে~$x$ পাওয়া যায়।

\begin{thm}\ollabel{thm:refl}
  $M \bredpar M$।
\end{thm}

\begin{proof}
  অনুশীলনী।
\end{proof}

\begin{prob}
  \olref[lam][cr][pb]{thm:refl} প্রমাণ করো।
\end{prob}

\begin{defn}[$\beta$-পূর্ণ বিকাশ]\ollabel{defn:bcd}
  $M$-এর \emph{$\beta$-পূর্ণ বিকাশ}~$\bcd{M}$ আরোহমূলকভাবে নিচের মতো
  সংজ্ঞায়িত:
  \begin{align}
    \bcd{x} &= x \ollabel{defn:bcd1} \\
    \bcd{(\lambd[x][N])} &= \lambd[x][\bcd{N}] \ollabel{defn:bcd2}\\
    \bcd{(PQ)} &= \bcd{P}\bcd{Q} && \text{যদি $P$ কোনো $\lambd$-বিমূর্ত পদ না হয়}
    \ollabel{defn:bcd3} \\
    \bcd{((\lambd[x][N])Q)} &= \Subst{\bcd{N}}{\bcd{Q}}{x} \ollabel{defn:bcd4}
  \end{align}
\end{defn}

নাম থেকেই বোঝা যায়, কোনো পদের $\beta$-পূর্ণ বিকাশ হলো একটি “পূর্ণ
সমান্তরাল হ্রাস”। সমান্তরাল $\beta$-হ্রাসে কোনো রিডেক্সকে সংকুচিত না
করার সুযোগ থাকে, কিন্তু পূর্ণ বিকাশে সব কটিকেই সংকুচিত করতে হয়। তাই
$(\lambd[f][fx])(\lambd[y][y])$-এর পূর্ণ বিকাশ পদটি নিজে নয়, বরং
$(\lambd[y][y])x$।

\begin{editorial}
  এই সংজ্ঞার একটি সমস্যা আছে: $\lambd$-পদের উপর অপেক্ষককে
  পুনরাবৃত্তিমূলকভাবে কীভাবে সংজ্ঞায়িত করতে হয়, তা এখনও প্রবর্তন করা
  হয়নি। ভবিষ্যতে এটি ঠিক করা হবে।
\end{editorial}

\begin{lem}\ollabel{lem:comp}
  $M \bredpar M'$ এবং $R \bredpar R'$ হলে $\Subst{M}{R}{y}
  \bredpar \Subst{M'}{R'}{y}$।
\end{lem}

\begin{proof}
  $M \bredpar M'$-এর !!{derivation}-এর উপর আরোহ প্রয়োগ করি।
  \begin{enumerate}
    \item শেষ ধাপটি \olref{defn:bredpar1}: অনুশীলনী।
    \item শেষ ধাপটি \olref{defn:bredpar2}: তখন $M$ হলো
      $\lambd[x][N]$ এবং $M'$ হলো $\lambd[x][N']$,
      যেখানে $N \bredpar N'$। আমাদের প্রমাণ করতে হবে
      $\Subst{(\lambd[x][N])}{R}{y} \bredpar
      \Subst{(\lambd[x][N'])}{R'}{y}$; অর্থাৎ,
      $\lambd[x][\Subst{N}{R}{y}] \bredpar
      \lambd[x][\Subst{N'}{R'}{y}]$। এটি \olref{defn:bredpar2} ও
      আরোহ-অনুমান থেকে অবিলম্বে অনুসৃত।
% BN-SRC-294 TARGET-CORRECTION-BEGIN
\par\smallskip\noindent\textnormal{\small\textbf{উৎস-সংশোধন (BN-SRC-294):} হিমায়িত বিমূর্তন-ক্ষেত্রের “অর্থাৎ” সম্বন্ধের ডান পাশে $\Subst{N'}{R}{y}$ লেখা, যদিও উপপাদ্যের সিদ্ধান্তে $\Subst{N'}{R'}{y}$ এবং আরোহ-অনুমানও $R\bredpar R'$ ব্যবহার করে ঠিক সেই ফল দেয়। বাংলা গণনায় হারিয়ে যাওয়া প্রাইমটি পুনরুদ্ধার করা হয়েছে।} % BN-SRC-294 TARGET-CORRECTION-END
    \item শেষ ধাপটি \olref{defn:bredpar3}: অনুশীলনী।
    \item শেষ ধাপটি \olref{defn:bredpar4}: $M$ হলো
      $(\lambd[x][N])Q$ এবং $M'$ হলো $\Subst{N'}{Q'}{x}$। আমাদের
      প্রমাণ করতে হবে $\Subst{((\lambd[x][N])Q)}{R}{y} \bredpar
      \Subst{\Subst{N'}{Q'}{x}}{R'}{y}$; অর্থাৎ,
      $(\lambd[x][\Subst{N}{R}{y}])\Subst{Q}{R}{y} \bredpar
      \Subst{\Subst{N'}{R'}{y}}{\Subst{Q'}{R'}{y}}{x}$। এটি
      \olref{defn:bredpar4} ও আরোহ-অনুমান থেকে অনুসৃত।
  \end{enumerate}
\end{proof}

\begin{prob}
  \olref[lam][cr][pb]{lem:comp}-এর প্রমাণ সম্পূর্ণ করো।
\end{prob}

\begin{lem}\ollabel{lem:cont}
  $M \bredpar M'$ হলে $M' \bredpar \bcd{M}$।
\end{lem}
\begin{proof}
  $M \bredpar M'$-এর !!{derivation}-এর উপর আরোহ প্রয়োগ করি।
  \begin{enumerate}
    \item শেষ বিধিটি \olref{defn:bredpar1}: অনুশীলনী।
    \item শেষ বিধিটি \olref{defn:bredpar2}: $M$ হলো
      $\lambd[x][N]$ এবং $M'$ হলো $\lambd[x][N']$, যেখানে
      $N \bredpar N'$। দেখাতে হবে $\lambd[x][N'] \bredpar
      \bcd{(\lambd[x][N])}$; অর্থাৎ, \olref{defn:bcd2} অনুসারে
      $\lambd[x][N'] \bredpar \lambd[x][\bcd{N}]$।
      \olref{defn:bredpar2} ও আরোহ-অনুমান থেকে এটি অনুসৃত।
    \item শেষ বিধিটি \olref{defn:bredpar3}: $M$ হলো $PQ$ এবং $M'$
      হলো $P'Q'$; এখানে কোনো $P$, $Q$, $P'$ ও~$Q'$-এর জন্য
      $P \bredpar P'$ এবং $Q \bredpar Q'$। আরোহ-অনুমান থেকে পাই
      $P' \bredpar \bcd{P}$ এবং $Q' \bredpar \bcd{Q}$।
      \begin{enumerate}
        \item কোনো $x$ ও~$N$-এর জন্য $P$ যদি $\lambd[x][N]$ হয়,
          তবে কোনো $N'$-এর জন্য $P'$-কে $\lambd[x][N']$ হতে হবে,
          যেখানে $N \bredpar N'$। আরোহ-অনুমান থেকে
          $N' \bredpar \bcd{N}$ এবং $Q' \bredpar \bcd{Q}$। তখন
          \olref{defn:bredpar4} অনুসারে $(\lambd[x][N'])Q' \bredpar
          \Subst{\bcd{N}}{\bcd{Q}}{x}$।
        \item $P$ কোনো $\lambd$-বিমূর্ত পদ না হলে
          \olref{defn:bredpar3} অনুসারে $P'Q' \bredpar
          \bcd{P}\bcd{Q}$; আর \olref{defn:bcd3} অনুসারে ডান পাশটি
          হলো $\bcd{PQ}$।
      \end{enumerate}
    \item শেষ বিধিটি \olref{defn:bredpar4}: $M$ হলো
      $(\lambd[x][N])Q$ এবং $M'$ হলো $\Subst{N'}{Q'}{x}$; এখানে
      কোনো $x$, $N$, $Q$, $N'$ ও~$Q'$-এর জন্য $N \bredpar N'$ এবং
      $Q \bredpar Q'$। আরোহ-অনুমান থেকে জানি
      $N' \bredpar \bcd{N}$ এবং $Q' \bredpar \bcd{Q}$।
      \olref{lem:comp} অনুসারে
      $\Subst{N'}{Q'}{x} \bredpar \Subst{\bcd{N}}{\bcd{Q}}{x}$;
      এখানে ডান পাশটি ঠিক $\bcd{((\lambd[x][N])Q)}$।
  \end{enumerate}
\end{proof}

\begin{prob}
  \olref[lam][cr][pb]{lem:cont}-এর প্রমাণ সম্পূর্ণ করো।
\end{prob}

\begin{thm}\ollabel{thm:cr}
  $\bredpar$ চার্চ--রসার ধর্ম পূরণ করে।
\end{thm}

\begin{proof}
  \olref{lem:cont} থেকে অবিলম্বে।
\end{proof}

\end{document}
